% SOURCE_STATEMENT: "There has been considerable progress (Kuijlaars,...) in understanding basic statistics such as the correlation functions for the 2-matrix random matrix model, and also matrix models with a source. The key element in these developments has been the successful extension by Kuijlaars et al of the Riemann-Hilbert/steepest descent method to $3\\times3$ Riemann-Hilbert problems. So far only the simplest situations have been considered. In order to consider the generic situation, one must, in particular, extend the Riemann-Hilbert/steepest descent method to $n\\times n$ Riemann-Hilbert problems. This is a challenging problem which would have important implications, not only for random matrix models, but also for problems in other areas, such as Pad\\'e-Hermite approximations and irrationality questions for distinguished real numbers, and multi-orthogonal polynomials."
% Preparation and final self-audit: CONTRIBUTING.md prompts C and D; source/audit artifacts retained in the local physics_statements_2026_10_11 batch.
% !TeX program = lualatex
\documentclass[11pt]{article}
\usepackage[margin=25mm]{geometry}
\usepackage{fontspec}
\setmainfont{Latin Modern Roman}
\usepackage{amsmath,amssymb,mathtools}
\usepackage{unicode-math}
\setmathfont{Latin Modern Math}
\usepackage{xurl,hyperref}
\hypersetup{colorlinks=true,urlcolor=blue}
\setcounter{secnumdepth}{0}
\setlength{\parindent}{0pt}
\setlength{\parskip}{5pt}
\setlength{\emergencystretch}{3em}
\newcommand{\sref}[2]{\hyperlink{source:#1:#2}{[S#2]}}
\newcommand{\cataloguescope}[1]{\paragraph{Recorded target.}#1}
\newcommand{\R}{\mathbb R}
\newcommand{\C}{\mathbb C}
\newcommand{\N}{\mathbb N}
\newcommand{\Z}{\mathbb Z}
\newcommand{\Q}{\mathbb Q}
\newcommand{\F}{\mathbb F}
\newcommand{\eps}{\varepsilon}
\begin{document}
\section{Multi-matrix models and models with an external field}\label{statement}

\subsection{Definitions and mathematical statement}

\paragraph{Shared notation.}$\mathbb N=\{1,2,\ldots\}$; the standard number fields have their usual meanings. All additional conventions are stated locally.


\paragraph{Definitions and notation.}A multi-matrix model is a probability law on several interacting random matrices; an external-source model couples a random matrix to a fixed matrix. Matrix-valued Riemann--Hilbert problems encode the relevant multiple orthogonal polynomials. The matrix size of the Riemann--Hilbert problem need not equal the size of the sampled random matrices.

\paragraph{Catalogue wording retained for traceability.}
There has been considerable progress (Kuijlaars,...) in understanding basic statistics such as the correlation functions for the 2-matrix random matrix model, and also matrix models with a source. The key element in these developments has been the successful extension by Kuijlaars et al of the Riemann-Hilbert/steepest descent method to $3\times3$ Riemann-Hilbert problems. So far only the simplest situations have been considered. In order to consider the generic situation, one must, in particular, extend the Riemann-Hilbert/steepest descent method to $n\times n$ Riemann-Hilbert problems. This is a challenging problem which would have important implications, not only for random matrix models, but also for problems in other areas, such as Pad\'e-Hermite approximations and irrationality questions for distinguished real numbers, and multi-orthogonal polynomials.
\subsection{Short English statement}
Extend Riemann--Hilbert steepest descent from special low-dimensional cases to general multi-matrix and external-source models.
\subsection{Sources}
\begin{itemize}
\item[\hypertarget{source:statement:1}{S1}] ULAM.AI and the UnsolvedMath Contributors, UnsolvedMath, catalogue entry. Catalogue curation under CC BY 4.0; definitions and English exposition added locally.
\url{https://huggingface.co/datasets/ulamai/UnsolvedMath}.
\item[\hypertarget{source:statement:2}{S2}] Percy Deift, Some open problems in random matrix theory and the theory of integrable systems (2007), arXiv:0712.0849, Problem 15.
\url{https://arxiv.org/abs/0712.0849}.
\end{itemize}
\label{end:statement}
\end{document}
